\begin{frame}[t]{跨模态 CPT：蛋白扰动评估及其边界}
\fontsize{8.8}{12}\selectfont\renewcommand{\arraystretch}{1.14}
\begin{tabularx}{\textwidth}{@{}P{2.9cm}Y Y@{}}
\toprule\textbf{评估数据}&\textbf{PXD014791}&\textbf{ProTargetMiner}\\\midrule
背景 / 类型 & 原代人心肌细胞，4 donors；LFQ & A549、MCF7、RKO；TMT 绝对 reporter 强度\\
药物与样本 & 主表 21 药；237 处理 + 71 对照 & 56 药；237 处理 + 35 对照\\
评估单位 & 60 个 drug--context 对，报告其中 58 对 & 73 个 drug--context 对，报告其中 61 对\\
Top-100 DE Pearson & 基础 70M：0.542；+蛋白 CPT：0.611 & 基础 70M：0.746；+蛋白 CPT：0.813\\
\bottomrule
\end{tabularx}
\vspace{.20cm}
\subhead{表示 + 近邻评估}
留一 drug--context 对；提取蛋白谱嵌入，按余弦相似度使用 $k=5$ 的训练对响应加权预测。
\remark{正文描述嵌入全部可用谱并按 drug--context 聚合；查询是否需要处理后谱尚不明确，不能直接等同 PT/G2CP 的“药物 + 对照 $\rightarrow$ 响应”。}
\source{[25] §2.3、Table 1、Fig. 3、SI。主表心肌集 21 药，而 SI 写 22 药中留出 21、保留 Trastuzumab；药物数口径需核对。论文报告增益，不代表本项目已复现。}
\end{frame}
